% The ablation subsection of point_defects_report.tex. It is read only when
% pd_table_ablations.tex exists (written by make_report_numbers.py once all runs are there).
\subsection{Ablations of the finder}
\label{sec:ablations}

% written by make_report_numbers.py
\begin{table}[t]
\centering
\caption{Ablations of the finder: defects found (\%) on the real test defects after 2,000 updates, one setting changed per run. Size in millions of parameters. Brackets: the range that holds the true rate with 95\% confidence. $p$: the probability that the difference from the run it is compared with is chance; that run is the reference, or, for the two runs with fewer layers, the run with random starting weights and 4 layers.}
\label{tab:ablations}
\footnotesize\setlength{\tabcolsep}{3.5pt}
\begin{tabularx}{\linewidth}{@{}>{\raggedright\arraybackslash}Xcccccll@{}}
\toprule
Run & Size & Vacancy & Interst. & Antisite & Frenkel & All & $p$ \\
\midrule
Reference: starts from the denoiser, 4 layers, class weights & 0.52 & 64 & 29.5 & 42.5 & 42 & 44.7 {\scriptsize[40, 49]} & -- \\
\addlinespace
Random starting weights & 0.52 & 17 & 27.5 & 26.5 & 24 & 23.7 {\scriptsize[20, 27]} & $<0.001$ \\
Random starting weights, 2 layers (reach $5\sigma$) & 0.29 & 15.5 & 27.5 & 25 & 24.7 & 23.1 {\scriptsize[20, 27]} & 0.46 \\
Random starting weights, 1 layer (reach $2.5\sigma$) & 0.17 & 13.5 & 23 & 15.5 & 17.3 & 17.3 {\scriptsize[14, 20]} & $<0.001$ \\
\addlinespace
No class weights & 0.52 & 61 & 31.5 & 41 & 36.7 & 42.9 {\scriptsize[39, 47]} & 0.12 \\
Trained without the Frenkel pairs that undid themselves & 0.52 & 64 & 29.5 & 38 & 35.3 & 42.1 {\scriptsize[38, 46]} & 0.027 \\
\bottomrule
\end{tabularx}
\end{table}


An ablation changes or removes one part of a method to see how much that part matters. Table~\ref{tab:ablations} does this for the finder, one setting per run. Six full training runs were beyond our computing budget, so each run has only \pd{abl-updates} updates, where the finder of the main results has 60,000. The rates are therefore lower than in Table~\ref{tab:found}, and the runs should be compared with each other only. There is one run per setting, so a difference of two or three points can also come from the randomness of training. The short runs are scored with their final weights, not with the moving average used for the long runs: an average over so few updates still contains 13\% of the starting weights, which would count against a run that starts from random weights.

\textbf{Starting from the denoiser matters most.} After the same number of updates, a finder that starts from random weights finds \pd{abl-scratch-all}\% of the defects, and the reference, which starts from the denoiser, \pd{abl-reference-all}\% (\pd{abl-p-scratch}). The denoiser was trained on defect-free glasses without any labels and already knows what normal glass looks like; the finder then only has to learn what a defect changes. The gain is largest for vacancies (\pd{abl-scratch-vacancy}\% against \pd{abl-reference-vacancy}\%), the defect that the denoiser alone already found best (Section~\ref{sec:iterations}). There is none for interstitials (\pd{abl-scratch-interstitial}\% against \pd{abl-reference-interstitial}\%), where the denoiser alone was poor. These short runs cannot say whether a network that starts from random weights would catch up with much longer training.

\textbf{The finder has to see beyond the nearest neighbours.} Each layer adds $2.5\sigma$ to the distance over which the label of an atom can depend on the structure. From random weights, one layer finds \pd{abl-layers1-all}\% of the defects, two layers \pd{abl-layers2-all}\% and four layers \pd{abl-scratch-all}\%. One layer is clearly worse (\pd{abl-p-layers1}); two and four do equally well at this budget (\pd{abl-p-layers2}). This agrees with the picture of a relaxed defect as a pattern over a few tens of atoms (Table~\ref{tab:physics}). Whether the full reach of four layers pays off cannot be seen in runs this short.

\textbf{Class weights change the labels, not the ranking.} Without the class weights the finder finds \pd{abl-no-weights-all}\% of the defects, and with them \pd{abl-reference-all}\%; the difference could be chance (\pd{abl-p-no-weights}). ``Found'' only asks which atom is ranked first, and the weights do not change that. They change how readily the network gives an atom a defect label. Trained without them, the network plays safe and calls nearly every atom intact, which is right 98\% of the time: on validation glasses it labels \pd{abl-label-no-weights-vacancy}\% of the atoms next to a vacancy as such, and no extra atom at all. With the weights these rates are \pd{abl-label-reference-vacancy}\% and \pd{abl-label-reference-interstitial}\%.

\textbf{Cleaning the training set did not help.} We expected the finder to improve when the Frenkel pairs that had undone themselves are taken out of the training set, since they carry defect labels but hold no defect. In this short run it did not: \pd{abl-real-only-frenkel}\% of the Frenkel pairs are found against \pd{abl-reference-frenkel}\%, and \pd{abl-real-only-all}\% against \pd{abl-reference-all}\% of all defects (\pd{abl-p-real-only}). We do not read this as a real loss. The $p$ value compares these two trained networks only; it says nothing about how much a result changes when the same training is simply repeated. Scored with the moving average of their weights, the two runs are equal (\pd{abl-avg-real-only}\% and \pd{abl-avg-reference}\%). What the run does show is that taking out the wrong labels brings no gain, at least early in training.
